\documentclass[11pt,a4paper]{article}
\usepackage[margin=2.3cm]{geometry}
\usepackage{amsmath,amssymb}
\usepackage{xcolor}
\usepackage{enumitem}
\usepackage{booktabs}
\usepackage{hyperref}
\usepackage{parskip}
\newcommand{\rev}[1]{\par\noindent\textbf{Reviewer comment.} \emph{#1}}
\newcommand{\resp}{\par\noindent\textbf{Response.} }
\newcommand{\chg}{\par\noindent\textbf{Changes.} }
\setlist{itemsep=2pt,topsep=2pt}
\definecolor{boxc}{gray}{0.95}
\begin{document}

\begin{center}
{\Large\bfseries Response to the Editor and Reviewers}\\[4pt]
Manuscript: \emph{Anytime-Valid Sequential Testing of Group Synergy: Interventional and Predictive Formulations, with an Application to Daily Financial Markets}\\
(previously: \emph{Anytime-Valid Nonparametric Testing of Synergistic (Group) Causality, with an Application Establishing Its Absence in Daily Financial Markets})\\
Journal: Discover Applied Mathematics\quad Submission ID: 7d12c900-4ad8-499d-aff6-602fc39f5a1e\quad Version 1.0 $\rightarrow$ revision
\end{center}

Dear Dr Wang and Ms Shinde,

Thank you for the opportunity to revise the manuscript and thank you to both reviewers for reports that were detailed, correct on the central points, and useful. We have revised the paper substantially. We respond to every comment below; section, table and figure numbers refer to the revised manuscript. All new experiments are in the public repository under \texttt{sequential-synergy-causality/experiments/revision/} (results in \texttt{results/}).

\section*{Summary of what changed}

\begin{enumerate}
\item \textbf{The validity proof for the group/time-series test was wrong, and is replaced (Reviewer 2, \S1.1; Reviewer 1, point 4).} Both reviewers are right that clipping does not preserve the sign of the conditional mean. We give Reviewer~1's counterexample in the paper (Example~1, and a simulation in which the clipped test falsely rejects $97.8\%$ of the time). The new construction uses a truncated loss and the predictable bound $b_t=2c|p^{\mathrm{add}}_t-p^{\mathrm{joint}}_t|$ on the score gap; dividing by $b_t$ preserves the sign of the conditional mean exactly, so validity (Theorem~8) needs no clipping and no distributional assumption. This also answers the request for robust losses (Reviewer~1, point 6).
\item \textbf{The theory is split into two parts and the nulls are restated (Reviewer 2, structural revision, \S1.2).} Part~A (interventional, Section~4) uses a sequential null $\theta_t\le0$ for all $t$; Part~B (predictive, Section~5) states the null directly as a comparison of two learned forecasters, $\mathbb E[s_t\mid\mathcal F_{t-1}]\le0$, with an explicit filtration paragraph (Section~5.1). ``Nonparametric'' no longer appears as a claim about the time-series procedure.
\item \textbf{The isolation proposition was over-claimed and is corrected (Reviewer 1, point 1; Reviewer 2, \S1.4).} We prove what holds (Proposition~10), show that the four-model contrast of the first version is contaminated by redundancy and suppression (it can be negative \emph{or positive} under a purely additive mechanism), switch to the add-versus-joint contrast that the reviewer proposed, and give counterexamples to the previous claim (including one single-source, fully Gaussian example).
\item \textbf{The financial conclusion is rewritten and the search is fully specified (Reviewer 1, points 3, 8; Reviewer 2, \S\S1.10--1.12).} The title, abstract and text no longer claim to ``establish absence''. A formal search of $27{,}027$ hypotheses over five panels, three target types and four lag structures is reported with family-wise thresholds, e-BH counts and a calibration on real data (circularly shifted placebo targets). The result is narrower than the first version's: in daily and weekly returns and volatility there is no family-wise detection and the real evidence matches the placebo null, but the same placebo shows that the test is \emph{not} calibrated for binary stress indicators and hourly crypto series (it rejects a true null in $10$--$40\%$ of hypotheses there), so the many rejections in those families are not interpreted; and the time-uniform bound on average synergy is weak relative to plausible effect sizes (Section~6.5).
\item \textbf{Problems we found ourselves while re-running everything.} We report them plainly because they affect the first version's claims: (a) the volatility and tail scans of the first version used contemporaneous features although the text called them lagged; (b) ``more than $3700$ triples'' double-counted one study; (c) standardisation used full-sample moments (look-ahead); (d) the turbulence example, presented as a positive control, is not valid: both the original and the revised test ``detect'' on phase-randomised surrogates that are independent of the sources (Section~6.4), so we withdrew it; (e) the portfolio-variance target is a deterministic function of the sources (Reviewer~1, point~9, is right) and is now presented only as a check of the contemporaneous variant.
\item \textbf{Declarations section} with the three requested subheadings (Ethical approval, Consent to participate, Consent to publish), and the reference list corrected (the Jha et al.\ preprint was missing from the list; its entry is now complete).
\end{enumerate}

Two points where we did not follow a suggestion are flagged as such below (Choi references; Clark--West as a test).

\section*{Editor's comments}

\rev{Citations: directly relevant prior work should be acknowledged; additional suggested references are at the authors' discretion. Declarations: add a section with Ethical approval, Consent to participate, Consent to publish.}
\resp We added the suggested references that are directly relevant and that we could verify: Billio et al.\ (2012), Choe and Ramdas (2024), Clark and West (2007), Gr\"unwald, Henzi and Lardy (2024), Podkopaev et al.\ (2023), together with Pocock (1977), O'Brien and Fleming (1979), Lan and DeMets (1983) for the group-sequential comparator, Rahimi and Recht (2007) for the random-feature dictionary, and Azuma (1967) and Wilson (1927) for the technical tools. We were unable to locate the two 2026 Choi / Choi--Kim records in the indexing services available to us, so we did not cite them; their substantive points (lag orders beyond one day, regime and crisis dependence of nonlinear transmission) are addressed by the lag-$5$ and weekly structures and the crisis analysis (Section~6.5). If the editorial office can supply DOIs we will add them. The contribution statement now restricts the novelty claim (Section~2.6).
\chg New ``Declarations'' section before the appendix with separate subheadings \emph{Ethical approval}, \emph{Consent to participate} and \emph{Consent to publish}, plus funding, competing interests, data and code availability (with software versions) and author contributions.

\section*{Reviewer 1}

We thank Reviewer~1 for the careful reading, in particular for points 1, 2, 4 and 9, each of which exposed a real problem.

\rev{(1) Does the pairwise score isolate interaction? With $s_t=\ell^A+\ell^B-\ell^{AB}-\ell^{\mathrm{base}}$ the population quantity is negative under redundancy even for additive mechanisms; Proposition~7's proof needs orthogonality. State which score each scan uses; re-run the pairwise scans with $\ell^{\mathrm{add}}-\ell^{\mathrm{joint}}$. ``$f$ arbitrary smooth'' is too broad.}
\resp You are right, and the problem is slightly worse than described. In Proposition~10(b) we show $\mathrm{Syn}_4=\mathrm{Syn}_{\mathrm{aj}}-R$ where the redundancy term $R$ has \emph{either sign}: for jointly Gaussian $(A,B)$ with correlation $\rho$ and $Y=A+B+\varepsilon$, $\mathrm{Syn}_4=-2\rho(1+\rho)$; for $Y=A-B+\varepsilon$ it is $+2\rho(1-\rho)$, i.e.\ positive ``synergy'' under a linear additive mechanism, while $\mathrm{Syn}_{\mathrm{aj}}=0$ in both. Empirically the four-model test rejects an additive null in $17\%$ of replications under persistent sources and $16\%$ at a realised source correlation of $0.20$ (its raw mean score is negative in those cells, so this is a finite-sample violation of its prequential null rather than redundancy as such), while the add-versus-joint test rejects in at most $1\%$ in all twelve stress nulls (Table~3, Figure~2). We therefore use the add-versus-joint score everywhere (Algorithm~2, Theorem~8) and \emph{all} scans in the paper were re-run with it, with the four-model score kept in the result files for comparison. On breadth: Proposition~10(c)--(d) states the exact conditions (own-span, independence, or one-lag Gaussian) under which $\mathrm{Syn}_{\mathrm{aj}}=0$ for an additive mechanism and gives counterexamples, including your point that a cubic term with a common factor can produce positive synergy (our counterexamples are $B=A^2+\tau V$, $Y=A^3+\varepsilon$ with $\mathrm{Syn}_{\mathrm{aj}}=36/(6+\tau^2)$, and a fully Gaussian two-lag example); the claim ``$f$ arbitrary smooth'' is withdrawn. The test itself remains valid for its prequential null regardless; what depends on the assumptions is only the reading of a rejection as an $A\times B$ interaction (Remark~3, Section~5.3).
\chg Sections 3.2, 5.3; Proposition~10 and its proof (Appendix A.2); Table~3, Figure~2 (stress tests), Figure~4 (correlated sources, correlation $0.2$--$0.99$, $150$ replications per cell); all financial scans re-run with the add-versus-joint score (Section~6.5, Table~6).

\rev{(2) Out-of-sample penalty on nested models; sensitivity to the forgetting factor; gap between the calibrated true-null median ($0.00$) and the financial median ($-0.93$).}
\resp We agree that the estimation-noise drift makes negative evidence uninformative, and say so (Section~5.4). \emph{Clark--West:} we report the Clark--West-adjusted score $s_t+(p^{\mathrm{add}}_t-p^{\mathrm{joint}}_t)^2$ as a diagnostic of the noise floor ($0.03$--$0.13$ for forgetting factors $0.99$--$0.999$), but we do not use it as a test statistic. Its conditional mean is $2(\mu_t-a_t)(j_t-a_t)$, which is not guaranteed to be $\le0$, so an e-process built on it would not be valid (Appendix A.2). \emph{Sensitivity:} forgetting factors $0.98,0.99,0.995,0.999$ on a null DGP, a heavy-tailed GARCH null and $150$ real triples give false-rejection rates $\le0.010$ in every simulated null cell ($100$ replications; the largest real-data rate is $0.047$ at $0.98$) and median final $\log_{10}E$ between $-0.04$ and $-0.27$; $0.98$ makes the learners unstable (Section~6.5). \emph{The gap:} with the revised test the simulated null median is $-0.15$ and the financial median $-0.27$; the placebo (circularly shifted) median on real data is $-0.20$ to $-0.26$. So the large gap in the first version came largely from the original score's normalisation. The remaining drift is concentrated in the first third of the sample (median contribution $-0.17$ finance, $-0.11$ null), the learning phase of the larger model, as the penalty predicts; it is not driven by winning bets that later reversed.
\chg Sections 5.4, 6.5 (paragraphs \emph{Null calibration on the real data} and \emph{Forgetting factor and the Clark--West diagnostic}); Appendix A.2.

\rev{(3) Evidence of absence is too strong: an e-value below $1$ does not bound the size of remaining synergy. Provide a confidence sequence or an equivalence-type e-process, state the minimum detectable effect, and explain how the calibrated synergy arm was chosen. The word ``Establishing'' in the title is too strong.}
\resp We agree and removed the claim. Title, abstract and conclusion no longer assert absence; the finding is now stated as ``no statistically significant super-additive predictive group effect in the settings examined''. We added (i) the reverse-hypothesis e-process $H_\delta:\mathbb E[s_t\mid\mathcal F_{t-1}]\ge\delta$ you suggested (Proposition~11, Section~5.5), with the exact statement of what its rejection excludes (persistent synergy at every monitored time, not an average) and a caveat that it can win mechanically when the two forecasters nearly agree; and, because that process turned out to be dominated by this effect on real data (the median smallest excluded $\delta$ was the grid minimum of $0.0005$, so we do not use it for conclusions), (ii) a time-uniform upper confidence bound for the \emph{running average} conditional synergy, in (truncated) loss units, with Bonferroni over the hypotheses of a family (Proposition~12, Section~6.5, \emph{How much is excluded}, Appendix Table~7). The bound is honest but weak: for lag-1 and contemporaneous daily families the median hypothesis excludes an average synergy above $0.057$--$0.122$ loss units, the worst hypothesis of a family only above $0.11$--$2.15$, and lag-5 and weekly families are barely informative. It therefore does not exclude effects of the size one would expect in daily returns (about $1\%$ of the variance); the paper says so. We note also that the units are truncated-loss units, not untruncated $R^2$ (the observations beyond three standard deviations are $2\%$ of the sample but $34$--$63\%$ of the sum of squares). We also measured power on real noise by injecting a known lagged interaction into real triples (Section~6.4, Figure~5). The result is sobering and we report it: power is low (at an injected term carrying $40\%$ of the variance, power is $0.43$ at $1/\alpha$ and $0.23$--$0.28$ at the family-wise thresholds), and relaxing the truncation and winsorisation (we tried $c=6$ without winsorisation) does not help and inflates the real-noise null rejection rate. The calibrated synergy arm of the first version was an arbitrary design choice (bilinear strength $0.9$, a very large effect), and the corresponding table was removed in favour of the placebo calibration and the injection study.
\chg Title, abstract, Sections 5.5, 6.4, 6.5, 7; Proposition~11; Figure~5.

\rev{(4) Proposition~6: clipping does not preserve the sign of the conditional mean. Restate the null on the clipped score or truncate losses; Choe and Ramdas (2024) is a template.}
\resp Agreed, and you were right that the squared-error setting is where this bites. See the summary above: Example~1 reproduces your calculation, a simulation confirms the failure for a fixed scale ($0.978$ false-rejection), we note that the null stated on the clipped score is valid but not the scientific null, and we replace the construction by the truncated-loss score with predictable bound $b_t=2c|p^1_t-p^2_t|$ (Lemma~13, Theorem~8, Corollary~9). Because $|s_t|\le b_t$ holds for every outcome (the truncated squared loss is $2c$-Lipschitz), $u_t=s_t/b_t\in[-1,1]$ with $\mathbb E[u_t\mid\mathcal F_{t-1}]=\mathbb E[s_t\mid\mathcal F_{t-1}]/b_t$, so the sign is preserved exactly and Ville's inequality applies. We adopt the forecaster-comparison null of Choe and Ramdas and cite them; their sequential-forecaster test is a natural baseline for future comparison.
\chg Section 5.2, Theorem~8, Corollary~9, Example~1, Appendix A.2, Table~3.

\rev{(5) Technical details: constant after ``absorbing'' $(1-\gamma)^{-1}$; Proposition~5 is only almost sure after a random time (use Freedman/Azuma); theorems concern an oracle $\lambda$ while the algorithms use a plug-in bet (regret argument); Theorem~2 is stated for $\theta=0$ although the null is one-sided; under nonstationarity cell means may vary.}
\resp All four points are fixed. (a) With $|\lambda\psi|\le\bar\gamma<1$ the sharp constant is $C_{\bar\gamma}=(-\log(1-\bar\gamma)-\bar\gamma)/\bar\gamma^2\le1/(2(1-\gamma))$; the growth bound is $g_0=\lambda\delta-C_{\bar\gamma}\lambda^2c^2$, and the cap-based Kelly bet is $\lambda^\star=\delta(1-\gamma)/c^2$ (Lemma~4, Theorem~5). (b) Proposition~6 is now a non-asymptotic high-probability bound by Azuma--Hoeffding with explicit tail $\exp(-2(tg_0-\log(1/\alpha))^2/(t\ell_\lambda^2))$ using the range $\ell_\lambda$ of the bet actually placed, plus an explicit $t_\beta$. We also corrected an inaccuracy: almost-sure \emph{convergence} of $t^{-1}\log E_t$ needs ergodicity; the theorem is now a $\liminf$. (c) We did not manage to prove a regret bound for the specific truncated plug-in bet and say so. We instead prove a regret statement for a mixture over a fixed grid of bets (Proposition~7: $\log E^{\mathrm{mix}}\ge\max_j\log E(\lambda_j)-\log K$), and Table~2 separates what is proved (validity for any predictable bet; power for the oracle or grid-mixture bet) from what is only evaluated empirically (power of the plug-in bet). (d) Theorem~2 is restated with a conditional contrast $\theta_t$, the one-sided null $\theta_t\le0$ for all $t$, a supermartingale for $\lambda\ge0$, and the two-sided average valid only under $\theta_t=0$; adaptive designs are allowed.
\chg Sections 4.1--4.3, Theorems 2 and 5, Propositions 6 and 7, Lemma 3, Table~2, Appendix A.1.

\rev{(6) Preprocessing: full-sample standardisation creates look-ahead. Squared loss is dominated by a few crisis days; report a robust loss or winsorisation. Extend lag order beyond $p=1$ (e.g.\ $1$--$5$ or weekly aggregation).}
\resp Standardisation is now \emph{predictable} (running mean and standard deviation of the past with a $30$-observation warm-up) for every real series. The loss is the truncated squared loss $\min(e^2,c^2)$, which is the robust loss; features are winsorised at $\pm5$; a sensitivity analysis over $c\in\{2,3,5\}$ and with/without winsorisation ($200$ hypotheses from each of two families) changes the raw rejection count by at most two hypotheses and leaves the median final $\log_{10}E$ between $-0.15$ and $-0.33$ (Section~6.5). The lag structure is extended to lag order $5$ (daily panels) and to weekly aggregation (non-overlapping $5$-day sums, lag $1$), in addition to lag $1$ and the contemporaneous variant (Section~6.5, Table~6). We note that the added robustness has a cost in sensitivity to tail-concentrated interactions, which we measure (Section~6.4).
\chg Sections 5.1, 5.2, 6.5; Table~6; \texttt{experiments/revision/r17\_sensitivity\_fwpower.py}.

\rev{(7) Predictor class is restricted; nonlinear/regime-specific synergy may exist. Report crisis-by-crisis trajectories; are reported $\log_{10}E$ values final values or running maxima? Podkopaev et al.\ (2023) could serve as a nonlinear extension.}
\resp Both final values and running maxima are now reported for every hypothesis (the result files contain $\log_{10}E$ at the final time and at its maximum, with the argmax time), and Table~6 and Figure~6 use the running maximum, which is the quantity the corollary refers to; the median final value is quoted separately. The trajectories across whole families (Figure~8) show that the maximum evidence over hypotheses peaks at $\log_{10}E=1.6$ (stock volatility, January 2017) and $2.2$ (ETF returns, 5 February 2020, just before the COVID sell-off) and decays afterwards, far from the family-wise thresholds ($4.1$, $4.4$); within the 2018, 2020 and 2022 windows the median of the within-window maxima is negative. We also address the dictionary restriction directly: a random-Fourier-feature dictionary detects an interaction outside the polynomial span (power $1.00$ versus $0.03$) with unchanged Type-I (Section~6.2). We cite Podkopaev et al.\ as a kernelised route to a nonlinear extension (Sections 2.5, 7).
\chg Sections 6.2, 6.5; Figures 6--8.

\rev{(8) Power at the family-wise thresholds actually used; e-BH discoveries; distribution of running-maximum e-values; threshold behind the $0.007$ true-null rate.}
\resp Done. Table~6 and Figure~6 give raw and family-wise counts and the distribution of running maxima per family; e-BH discoveries at FDR $0.1$ and $0.05$ (using final e-values, which are the valid input) are in the result file \texttt{r5\_summary.json}. Power at the family-wise thresholds ($m=660$, $m=4400$) is reported on real noise (Section~6.4, Figure~5) and for the synthetic interaction DGP (Section~6.1): at interaction strength $0.2$ power falls from $0.83$ at $1/\alpha$ to $0.06$ and $0.03$ at the $m=660$ and $m=4400$ thresholds, at $0.3$ it is $0.99$, $0.77$ and $0.61$, and at $0.45$ it is $1.00$ throughout. The old $0.007$ was the rate at which $300$ simulated null series crossed $1/\alpha=20$; that table is removed.
\chg Sections 6.1, 6.4, 6.5; Table~6; Figures 5, 6.

\rev{(9) The portfolio realised-variance control: the target is a deterministic function of the pair; puzzle with near-collinear pairs and equity--bond sign. Suggest a semi-synthetic control.}
\resp You are right. We now state the target explicitly, $RV_t=(0.5r_A+0.5r_B)^2$, describe the experiment as a check that the contemporaneous variant can recover a known interaction (not as group causality), and explain both puzzles. Near-collinear equity pairs have the weakest signal because the cross product has $R^2=0.96$--$0.995$ on the own-feature block $(1,r_A,r_B,r_A^2,r_B^2)$ (cross-asset pairs: $0.03$--$0.45$), so the cross term is nearly absorbed by the additive model; equity--bond pairs with negative covariance are detected strongly because the test measures predictive gain from the cross term, which does not depend on the sign of the covariance. All $14$ pairs are detected at $1/\alpha$ and $13$ at the family-wise threshold. We also added the semi-synthetic control you suggested: a known lagged interaction injected into real residuals (Section~6.4, Figure~5).
\chg Section 6.4 (paragraphs \emph{Contemporaneous portfolio variance} and \emph{Semi-synthetic power on real noise}); Figure~5.

\rev{(10) Benchmark: use a group-sequential design with alpha spending as comparator; explain $0.031$ vs $0.016$; describe how SURD/PEID are applied to continuous VAR data; Jha's test is missing from the reference list; Gaussian PID synergy is a different estimand; quantify the pairwise layer (Billio et al.).}
\resp Group-sequential Pocock and O'Brien--Fleming designs with Lan--DeMets spending were added in both the interventional experiment ($20$ looks, $2000$ streams: Type-I $0.051$ for both, matching the design; detection times are comparable to ANTE's when the schedule is fixed in advance) and the predictive benchmark ($15$ looks; Table~4). The $0.031$ and $0.016$ referred to different designs (interventional, $2000$ streams of length $4000$; predictive VAR with $15$ looks, $250$ streams of length $3000$), a difference we now state and replace with intervals. The SURD and PEID settings are given in Section~6.3 (six equal-mass bins on the lagged triple; polynomial-lifted Gaussian MI under an independent maximum-entropy intervention) and all code is released. The distributional Granger test attributed to Jha appeared in the tables but was not cited in the text and its bibliography entry was incomplete; it is now cited and completed (Jha, 2026, arXiv:2606.22230). We rephrase the PID/SURD/PEID results as a difference in estimand and no longer call them spurious (Section~6.3). The pairwise layer is quantified with Granger-causality networks on the same panels (Section~6.5, last paragraph): after Bonferroni, $20\%$ of ordered ETF pairs and $78\%$ of volatility pairs are significant at lag $1$, with a rolling density up to $0.55$ around the 2020 crash.
\chg Sections 6.1, 6.3, 6.5; Table~4; bibliography.

\rev{(11) ``Predictive synergy'' and ``group causality'' are used interchangeably; clarify in title and abstract; consider anytime-valid conditional-independence tests (Gr\"unwald et al.\ 2024).}
\resp Terminology is now fixed: \emph{interaction}, \emph{synergy}, \emph{causal synergy} (interventional design only) and \emph{predictive group synergy} (time series) are defined in the introduction, and the time-series procedure is never called causal. The title no longer contains ``causality''. Gr\"unwald et al.\ is cited and named as the route to a causal reading and to conditioning on confounders (Sections 2.5, 7); the AIPW extension is presented as an extension, with a proof only for the known-design case (Section~4.4).
\chg Title, abstract, Sections 1, 3, 4.4, 7.

\rev{(12) Presentation: Fig.~2 monotonicity; ``more than three thousand'' vs ``more than $3700$''; Table~6 column ``Real synergy''; ``a Q1 test''; Fig.~5 has panels B and C but no A; repository URL lost underscores; VIX is not traded; last sentence of \S5.4 mentions biology.}
\resp All fixed. The power/latency figure is regenerated with $500$ replications and bands (Figure~3); the claims of monotonicity are now stated within intervals, with the original test's dips noted. The hypothesis count is now exact ($27{,}027$; Section~6.5) and the old figure's double counting is explained. The ``Real synergy'' table and the panel-labelled figure are removed. The informal phrase and the sentence on biology are deleted. The repository URL is typeset with \verb|\url|. A note on VIX is added to Table~5. We also thank the reviewer for the careful eye.

\section*{Reviewer 2}

We thank Reviewer~2 for a rigorous report. We follow its numbering.

\paragraph{1.1 Validity of Proposition~6 (old numbering) and the clipping argument.} Agreed; this was the most important problem and is resolved by option (c)/(d) of the report: the construction is replaced by a truncated-loss score with a predictable bound (Theorem~8) and the assumptions are explicit; the clipped procedure is recorded as valid only for its own clipped null or under conditional symmetry (Appendix A.2), with Example~1 and a simulation. Simulations are retained as a complement, not as a substitute (Table~3).

\paragraph{1.2 Fixed causal parameter versus sequential null.} Restated: Assumption~1 (known, possibly adaptive, sequential randomisation), $\theta_t=\mathbb E[\theta_t(X_t)\mid\mathcal F_{t-1}]$, $H_0^{\le}:\theta_t\le0\ \forall t$, $H_0^{=}:\theta_t=0\ \forall t$ (Section~3.1, Lemma~1, Theorem~2). The statement that ANTE ``requires neither stationarity nor correct model specification'' is limited to what is proved (the sequential null; the prequential null).

\paragraph{1.3 Meaning of ``nonparametric''.} The word is removed from the title and from the claims about the time-series procedure, which is described as a model-based prequential implementation of an anytime-valid testing framework; the distribution-free statement is limited to the interventional part under boundedness and centring. A nonlinear dictionary (random Fourier features) is evaluated to show the dependence on the predictor class (Section~6.2).

\paragraph{1.4 Proposition~7.} Corrected as described in our response to Reviewer~1, point~1 (Proposition~10): explicit assumptions, proofs, counterexamples; the strong ``provably isolates'' claim is withdrawn. The robustness simulations over source correlation requested in 2.3 are in Section~6.2.

\paragraph{1.5 Causal versus predictive synergy.} Terminology defined in Section~1 and kept throughout; ``synergistic Granger causality'' is replaced by ``predictive group synergy''. Limitations state that the predictive null tests super-additive predictability, not a structural intervention (Sections 3, 5.1, 7).

\paragraph{1.6 ``First'' and ``no prior method''.} Replaced by ``to the best of our knowledge, ANTE is the first framework we are aware of that combines'' (P1) anytime-valid inference for an interaction contrast and (P3) sequential inference for groups, with an interventional and a predictive instantiation; the five properties you list are separated in Table~1 (Section~2.6). The abstract no longer uses ``every existing test'' or ``uniquely''.

\paragraph{1.7 Benchmark does not support blanket superiority.} The text and table now describe the comparison as complementary: on detection AUC SURD and PEID are as accurate as the revised test ($1.00$ each; the original test $0.992$); ANTE's distinct property is validity at any stopping time. The capability matrix with a ``winner'' column and phrases such as ``the only method that qualifies'' were removed (Section~6.3, Table~4).

\paragraph{1.8 Isolation comparison.} PID/SURD/PEID are described as answering a different question (information that is only jointly available), not as producing spurious synergy; ANTE's contrast is described as designed to vanish under additive mechanisms, with the conditions of Proposition~10 (Section~6.3).

\paragraph{1.9 The contemporaneous financial example.} Renamed and explained: the target is a deterministic function of the sources; it checks that the contemporaneous variant recovers a known interaction; we distinguish contemporaneous statistical interaction, predictive group synergy and causal synergy (Section~6.4). The section title no longer says ``genuine contemporaneous synergy''.

\paragraph{1.10 The financial ``absence'' claim.} Replaced by your suggested form: \emph{no statistically significant super-additive predictive group effect was detected in the financial datasets, frequencies, targets, conditioning sets and horizons examined} (Sections 6.5, 6.6, 7), qualified by the power limitations we measured. Sentences such as ``daily markets carry no irreducible group synergy'' are removed.

\paragraph{1.11 Power.} The calibrated synergy arm is no longer used as evidence that the test is powerful against arbitrary financial synergy. We report power on real noise with a known injected interaction (low for tail-concentrated interactions), against several synthetic alternatives, and as a function of source correlation, and we state the consequence for the interpretation of non-rejection (Sections 6.1, 6.2, 6.4, 6.6).

\paragraph{1.12 Multiple testing and search space.} Section~6.5 answers each of your eight items: (1) $27{,}027$ hypotheses in $43$ families ($\approx26{,}900$ distinct after removing $132$ duplicated tuples), listed in Table~5/6 and in \texttt{r5\_search\_space.json}; (2) \emph{not} pre-registered: the first version's scans were added sequentially, and the present grid was fixed before it was run but designed after the earlier results were known (stated in the text); (3) e-values are aggregated by e-Bonferroni (running maximum, valid under arbitrary dependence) and e-BH on final e-values; (4) per-family thresholds $m/\alpha$ and a global threshold $\log_{10}(27{,}027/0.05)=5.73$; (5) a hypothesis family is a (panel, target, lag structure) setting; (6) the panels were chosen as standard sets by the author rather than by selecting on outcomes, but the first version's results informed which structures to add; (7) overlapping hypotheses are counted and handled by Bonferroni; (8) yes, they overlap in assets and in targets. The first version's count of ``$>3700$'' double-counted one study and is corrected.

\paragraph{1.13 Forgetting factor and the filtration.} Section~5.1 (paragraph \emph{Filtration and predictability}) states that $p^k_t=\langle w^k_{t-1},\phi^k_t\rangle$ is $\mathcal F_{t-1}$-measurable (prediction precedes the update), and lists the other ingredients (standardisation, features, normaliser, bet, warm-up) and the requirement that all tuning constants be fixed in advance; Lemma~13 and Theorem~8 use it.

\paragraph{1.14 Consistency and the adaptive betting rule.} The text now distinguishes validity for any predictable bet, power for an oracle fixed bet or a fixed-grid mixture, and the behaviour of the plug-in bet, which is evaluated empirically only (Section~4.3, Table~2).

\paragraph{1.15 Detection-time bound.} Proposition~6 is a non-asymptotic \emph{high-probability} bound with explicit constants (not an expectation bound, not asymptotic), see also our reply to Reviewer~1, point~5.

\paragraph{1.16 AIPW.} Presented as an extension; the finite-sample guarantee is proved only for a known design and past-fitted outcome models (Section~4.4); the estimated-propensity case is explicitly not covered.

\paragraph{2.1 Type-I experiment.} Monte Carlo intervals added everywhere; the repeated Wald procedure is labelled as an illustration of the optional-stopping failure and not as a competing anytime-valid method; a group-sequential comparator is added (Section~6.1).

\paragraph{2.2 Benchmark sample sizes and tuning.} Section~6.3 and \texttt{results/r9\_benchmark\_ci.json} list sample sizes, replications, bins and estimator settings, the per-replication seeds, and the fact that no method is tuned separately.

\paragraph{2.3 Synthetic DGPs with correlated sources.} Added: source correlation from $0.2$ to $0.99$ with a true additive mechanism (Section~6.2, Figure~4).

\paragraph{2.4 Nonlinear alternatives.} Added: interactions inside and outside the polynomial span, and a random-feature dictionary (Section~6.2).

\paragraph{2.5 Physics example.} Re-examined, and withdrawn as evidence: on phase-randomised surrogates that are independent of the sources, both the original and the revised test falsely reject in $20/20$ cases, so the turbulence signals (smooth, nearly deterministic) cannot validate either test (Section~6.4). The reader should treat the first version's turbulence result as unsupported. The $k$-way test is validated on synthetic three-way mechanisms.

\paragraph{3 Financial application.} Conclusions are limited to the searched space, as proposed; the search is described precisely (Section~6.5).

\paragraph{4 Minor comments.} (1) abstract softened; (2) title changed; (3) terminology defined (Section~1); (4) notation made consistent ($\theta_t$; $\mathrm{Syn}_4$, $\mathrm{Syn}_{\mathrm{aj}}$); (5) filtration stated (Section~5.1); (6) $\hat\sigma_{t-1}$ no longer exists: the normaliser is $b_t$ with the rule $u_t=0$ when $b_t\le\varepsilon$ (Section~5.2); (7) the plug-in Kelly bet uses $\lambda_t=0$ when $\hat v_{t-1}=0$ (Section~4.2); (8) the warm-up and caps are part of the definition of the e-process (Sections 4.2, 5.1); (9) the one-sided process is used for the scientific question and the two-sided one is only for antagonism (Remark~1); (10) XOR: a negative additive interaction ($\theta=-2$), which differs from the information-theoretic name ``pure synergy'' (Section~3.1); (11) ``inferior'' is not used; (12) ``collapses'' replaced; (13) Monte Carlo intervals for all rates and AUCs; (14) hypotheses per family in Table~5/6; (15) not pre-specified, disclosed; (16) ``exhaustively'' removed; (17) software versions and environment added (Declarations); (18) Table~2 separates guarantees from empirical properties; (19) the conclusion separates the methodological result from the empirical finding.

\paragraph{5 Suggested structural revision.} Adopted: Part~A (Section~4) and Part~B (Section~5), with the predictive null stated directly as $\mathbb E[s_t\mid\mathcal F_{t-1}]\le0$ and the RLS implementation presented as one forecasting instantiation.

\paragraph{6 Overall evaluation, items 1--8.} (1) proof repaired; (2) assumptions for the isolation result stated and proved; (3) causal and predictive synergy separated; (4) ``nonparametric'' scope clarified; (5) proved properties separated from those of the RLS implementation (Table~2); (6) the multiple-testing procedure and search space are accounted for; (7) novelty and ``absence'' claims moderated; (8) simulations under correlated sources and nonlinear alternatives added.

\bigskip
Sincerely,\\
Quang-Vinh Dang\\
British University Vietnam, Hung Yen, Vietnam\quad vinh.dq4@buv.edu.vn\quad ORCID 0000-0002-3877-8024

\end{document}
